%!TEX root = ../template.tex
%%%%%%%%%%%%%%%%%%%%%%%%%%%%%%%%%%%%%%%%%%%%%%%%%%%%%%%%%%%%%%%%%%%%
%% chapter4.tex
%% NOVA thesis document file
%%
%% Results
%%%%%%%%%%%%%%%%%%%%%%%%%%%%%%%%%%%%%%%%%%%%%%%%%%%%%%%%%%%%%%%%%%%%

\typeout{NT FILE chapter4.tex}%

\chapter{Statistics and Results}
\label{chap:results}

This chapter reports what the learned observation-space adversary of
Section~\ref{subsec:learned_adversary} actually achieves against the victim family of
Section~\ref{sec:victim_architectures}. The question inherited from the earlier
stages of this research line is narrow and comparative. The FGSM study found that
perturbing every agent at every step altered a substantial fraction of routing
decisions while extracting almost none of the available damage. The present work asks
whether that robustness survives a stronger attacker with two mechanisms. Inspired by
PGD, it divides its effort across time: rather than spending a full-strength
perturbation at every step, it acts on only a fraction of the steps, choosing
\emph{when} to spend a limited budget. Inspired by SAJA, it coordinates the perturbation
across agents, choosing the compromised agents' perturbations together rather than one
agent at a time.

Both mechanisms are evaluated. The independent scope perturbs each agent separately, so
it isolates the first mechanism, dividing the attack across time. The coordinated scope
adds coordination, and comparing the two scopes on identical episodes isolates what
coordination contributes. Results are reported under two domain clamps, which decide
which observation features are re-projected onto $[0,1]$ after the perturbation.
Section~\ref{sec:results_parity_clamp} reports the independent scope under the
\emph{FGSM-parity} clamp, which reproduces the admissible set of the FGSM study, so
that the attack is measured on the same terms as the baseline.
Section~\ref{sec:results_full_clamp} reports both scopes under the stricter
\emph{full} clamp, and Section~\ref{sec:results_clamp_effect} compares the two clamps
on paired episodes.

The results are organised around three measurements. The first is how the attack's
effect varies with the timing budget, that is, with the fraction of timesteps the
adversary is permitted to act on; this tests the central claim of
Section~\ref{subsec:timed_attack}, that concentrating a fixed budget on
high-utilisation moments buys more damage than spreading it uniformly. The second is
the relationship between decisions changed and packets lost, which tests whether the
flip rate reported by the earlier study is a meaningful proxy for damage at all. The
third is the effect of coordination, measured as the paired difference between the
coordinated and independent attacks.

% ─────────────────────────────────────────────────────────────────────────────
\section{Reported Statistic and Experimental Protocol}
\label{sec:results_protocol}

Every figure and table in this chapter reports the \emph{adversarial gap} of
Equation~\eqref{eq:adversarial_gap}: the paired per-episode difference in end-to-end
packet delivery ratio (PDR) between the random control arm and the attack arm,
expressed in percentage points. The raw drop from the clean baseline is deliberately
not used as the headline number. A perturbation of magnitude $\varepsilon$ applied to
an observation will disturb the controller whether or not its direction was chosen
adversarially, so the clean-to-attack difference conflates the damage attributable to
the attacker's optimisation with the damage any same-sized noise would have caused.
The adversarial gap isolates the former, and it is the only quantity that speaks to
whether learning the perturbation direction was worth anything.

Three arms are run for each configuration: a clean arm with no perturbation, a
random control drawing uniformly from the same $\varepsilon$-ball under the same
admissible set and the same timing gate, and the attack arm. The three arms are
\emph{paired}: each of the 30 evaluation episodes uses an identical traffic
sequence and identical link-failure seed across all arms, so the per-episode
differences are matched and the reported 95\% confidence intervals are computed on
those paired differences. Without pairing, the between-episode variance in offered
traffic would dominate the effect being measured and the intervals would be
meaningless.

The victim is frozen throughout, as established in
Section~\ref{subsec:threat_model_pgd}: neither actor nor critic parameters receive
gradient updates during any attack or evaluation, so every number below describes an
attack on a fixed controller rather than a training-time interaction. Evaluation
uses $\varepsilon = 0.30$ under the $L_\infty$ norm, 30 paired episodes of 256
timesteps at an offered load factor of 2.0, with no injected link failures and a
fixed traffic seed. Three sets of adversaries, one per victim variant in each set,
were trained as described in Section~\ref{subsec:learned_adversary}: an independent
and a coordinated set under the full domain clamp, and an independent set under the
FGSM-parity clamp. An adversary is always evaluated under the clamp it was trained
with. Each was trained for 300 episodes with the timing gate set to $b = 0.25$, the
loosest budget in the sweep, so that the tighter evaluation budgets act on a subset of
the state distribution seen during training.

A result is treated as evidence of a real effect only when its 95\% confidence
interval excludes zero; such cells are marked in the tables and drawn with filled
markers in the figures. A confidence interval spanning zero is reported as a null,
not as proof of robustness, since each run has a finite resolution floor that
depends on the variant's paired variance. For the tightest variants this floor is
approximately $0.1$~pp, meaning a true effect smaller than that would be invisible
at $n = 30$ and should not be claimed either way.

% ─────────────────────────────────────────────────────────────────────────────
\section{Results Under the FGSM-Parity Domain Clamp}
\label{sec:results_parity_clamp}

Every result in this section was obtained under the \emph{FGSM-parity domain clamp}:
after the $\ell_\infty$ projection, only the normalised bandwidth features of each
agent's observation are re-projected onto $[0,1]$, and every other feature may move
anywhere within the $\varepsilon$-ball. This is the domain constraint $\mathcal{D}$ of
Equation~\eqref{eq:admissible_set}, the one used by the FGSM study of
Section~\ref{subsec:base_fgsm}, so the attack and its random control draw from exactly
the admissible set within which the baseline perturbed. A separate set of seven
independent-scope adversaries was trained for this clamp, which is recorded in each of
the 35 evaluation records. The clean arm depends on neither the adversary nor the
clamp, and its per-episode delivery series is bit-identical to that of the full-clamp
evaluation in all 35 configurations. The coordinated scope has not been evaluated under
this clamp.

This clamp is the looser of the two: every perturbation admissible under the full clamp
of Section~\ref{sec:results_full_clamp} is also admissible here. That does not make the
damage measured here an upper bound on the damage measured there. Each clamp has its
own trained adversary, and a learned adversary is not guaranteed to exploit a larger
set at least as well as a smaller one; Section~\ref{sec:results_clamp_effect} shows
that on several variants it does not.

% ─────────────────────────────────────────────────────────────────────────────
\subsection{Sensitivity to the timing budget}
\label{subsec:results_parity_timing}

\begin{figure}[htbp]
  \centering
  \includegraphics[width=\textwidth]{timing_sweep_gap_independent_parity.png}
  \caption{Independent scope, FGSM-parity domain clamp. Adversarial gap for all seven
           victim variants at $\varepsilon = 0.30$, $n = 30$ paired episodes.
           \textbf{Left:} gap against the requested timing budget, with the mean
           realised attack rate printed beneath each tick and the ungated run
           included as the 100\% reference; filled markers denote a 95\% confidence
           interval excluding zero, hollow markers a null. \textbf{Right:} the same
           gaps plotted against the fraction of routing decisions the attack
           actually changed, on a logarithmic axis, with each variant drawn as one
           path walking from $b = 0.10$ through to the ungated run. A thin line joins
           a label to its series where the label had to be moved clear of the data.}
  \label{fig:sweep_parity}
\end{figure}

The timing gate spends the budget it is given: the four requested budgets realise,
averaged over the seven variants, as $10.0\%$, $15.0\%$, $19.9\%$ and $24.9\%$ of
steps. Figure~\ref{fig:sweep_parity} and Table~\ref{tab:gap_sweep_parity} show a
different architectural pattern from the one found under the full clamp. A single
variant is clearly vulnerable. \texttt{LC-Simple} is attacked significantly at every
budget, its gap rising from $+0.40$~pp at $b = 0.10$ to $+1.45$~pp at $b = 0.25$ and
$+5.90$~pp ungated, the largest gap measured in this study. Ungated, its delivery ratio
falls to $71.8\%$, against $77.3\%$ clean and $77.7\%$ under the random control.
\texttt{CC-Duelling-GNN} separates from zero only ungated, at $+0.78$~pp, and
\texttt{CC-Duelling} does not separate at all, reaching $+0.59$~pp ungated with an
interval that spans zero. \texttt{CC-Simple}, \texttt{CC-Simple-GNN} and
\texttt{LC-Duelling-GNN} stay within $0.2$~pp of zero at every budget.

\begin{table}[htbp]
  \centering
  \caption{Independent scope, FGSM-parity domain clamp. Adversarial gap in percentage
           points by victim variant and requested timing budget, at
           $\varepsilon = 0.30$ and $n = 30$ paired episodes. Bold entries have a 95\%
           confidence interval excluding zero. Positive values indicate the attack
           cost the victim more than the random control of the same magnitude;
           negative values indicate the reverse.}
  \label{tab:gap_sweep_parity}
  \begin{tabular}{lrrrrr}
    \hline
    \textbf{Variant} & \textbf{0.10} & \textbf{0.15} & \textbf{0.20} & \textbf{0.25} & \textbf{ungated} \\
    \hline
    CC-Simple        & $-0.01$ & $+0.02$ & $-0.01$ & $+0.05$ & $-0.02$ \\
    CC-Duelling      & $-0.04$ & $+0.10$ & $+0.18$ & $+0.09$ & $+0.59$ \\
    CC-Simple-GNN    & $+0.00$ & $+0.00$ & $-0.01$ & $+0.00$ & $+0.04$ \\
    CC-Duelling-GNN  & $-0.17$ & $-0.24$ & $+0.11$ & $-0.15$ & $\mathbf{+0.78}$ \\
    LC-Simple        & $\mathbf{+0.40}$ & $\mathbf{+1.01}$ & $\mathbf{+1.27}$ & $\mathbf{+1.45}$ & $\mathbf{+5.90}$ \\
    LC-Duelling      & $\mathbf{-0.48}$ & $\mathbf{-0.71}$ & $\mathbf{-1.04}$ & $\mathbf{-0.83}$ & $\mathbf{-1.31}$ \\
    LC-Duelling-GNN  & $+0.00$ & $+0.01$ & $+0.01$ & $+0.00$ & $-0.18$ \\
    \hline
  \end{tabular}
\end{table}

\texttt{LC-Duelling} is an exception in the other direction. Its gap is significantly
\emph{negative} at every budget, from $-0.48$~pp at $b = 0.10$ to $-1.31$~pp ungated.
Ungated, the victim delivers $78.6\%$ of packets under attack, against $77.3\%$ under
the random control and $76.9\%$ with no perturbation at all, so the adversary trained
against this victim steers it to deliver more than it would unperturbed. This is not
evidence of robustness. It shows that the direction the adversary learned is
systematically wrong for this victim. The adversary maximises the per-step drop
fraction of Equation~\eqref{eq:la_reward} rather than the end-to-end delivery ratio on
which it is scored, and either the two objectives come apart on this victim or training
converged to a poor direction; these runs do not separate the two.

Under this clamp, concentrating the budget on critical states does not buy more damage
per attacked step. Normalising \texttt{LC-Simple}'s gap by the realised attack rate
gives $4.0$~pp per unit rate at $b = 0.10$, $6.7$ at $b = 0.15$, $6.4$ at $b = 0.20$
and $5.9$ at $b = 0.25$, against $5.9$ ungated. Beyond the tightest budget its damage
grows in proportion to the number of steps attacked, and the tightest budget is the
least efficient rather than the most. No other variant has a significant positive gap
at any gated budget, so none shows an efficiency gain either. The efficiency found under
the full clamp in Section~\ref{subsec:results_ind_timing} is therefore a property of
the adversaries trained there, not of dividing the attack across time as such. As under
the full clamp, the adversaries were trained with the gate at $b = 0.25$, so the
ungated column asks them to act on states they did not encounter during training. For
\texttt{LC-Simple} this costs nothing: its ungated gap is four times its gap at
$b = 0.25$, for four times the number of steps attacked.

% ─────────────────────────────────────────────────────────────────────────────
\subsection{Decision flips versus delivered damage}
\label{subsec:results_parity_flips}

The right panel of Figure~\ref{fig:sweep_parity} and
Table~\ref{tab:ungated_detail_parity} show the flip rate failing as a measure of damage
more starkly than under the full clamp. Ungated, \texttt{LC-Simple},
\texttt{CC-Duelling-GNN} and \texttt{LC-Duelling} change $15.3\%$, $13.3\%$ and
$13.2\%$ of routing decisions, within about two percentage points of one another, yet
their gaps are $+5.90$, $+0.78$ and $-1.31$~pp, a spread of more than seven percentage
points. The same holds at a gated budget: at $b = 0.25$, \texttt{LC-Simple} and
\texttt{CC-Duelling-GNN} both change $2.92\%$ of decisions, with gaps of $+1.45$ and
$-0.15$~pp. The random control again changes many decisions on its own, between $9.6\%$
and $10.7\%$ on the four variants with the highest flip rates.

\begin{table}[htbp]
  \centering
  \caption{Independent scope, FGSM-parity domain clamp. Ungated results at
           $\varepsilon = 0.30$, $n = 30$ paired episodes: adversarial gap with its
           95\% confidence interval, and the fraction of routing decisions changed by
           the attack and by the random control.}
  \label{tab:ungated_detail_parity}
  \begin{tabular}{lrcrr}
    \hline
    \textbf{Variant} & \textbf{Gap (pp)} & \textbf{95\% CI} & \textbf{Attack flips} & \textbf{Random flips} \\
    \hline
    CC-Simple        & $-0.02$ & $[-0.17, +0.14]$ & $8.68\%$  & $5.23\%$  \\
    CC-Duelling      & $+0.59$ & $[-0.08, +1.26]$ & $10.77\%$ & $9.62\%$  \\
    CC-Simple-GNN    & $+0.04$ & $[-0.02, +0.10]$ & $1.02\%$  & $0.27\%$  \\
    CC-Duelling-GNN  & $+0.78$ & $[+0.20, +1.36]$ & $13.34\%$ & $10.63\%$ \\
    LC-Simple        & $+5.90$ & $[+4.87, +6.93]$ & $15.32\%$ & $10.74\%$ \\
    LC-Duelling      & $-1.31$ & $[-1.94, -0.68]$ & $13.15\%$ & $9.87\%$  \\
    LC-Duelling-GNN  & $-0.18$ & $[-0.42, +0.06]$ & $3.11\%$  & $1.11\%$  \\
    \hline
  \end{tabular}
\end{table}

The two GNN variants that resist the attack show the same two behaviours as under the
full clamp (Section~\ref{subsec:results_ind_flips}). \texttt{CC-Simple-GNN} changes
only $1.02\%$ of decisions ungated, and at most $0.13\%$ at any gated budget, so the
perturbation rarely moves the selected path at all. \texttt{LC-Duelling-GNN} changes
$3.11\%$ ungated and still returns a gap of $-0.18$~pp: its routing does change, and
the network absorbs the change.

% ─────────────────────────────────────────────────────────────────────────────
\subsection{What triggered the attacks}
\label{subsec:results_parity_triggers}

These are the first runs in which the timing gate recorded why it acted. For every
attacked step it logs which case of Equation~\eqref{eq:timing_rule} admitted the step,
a score strictly above the threshold or a tie at the threshold admitted by rationing,
together with the congestion band the score fell in: saturated ($c_t \ge 1.00$), near
saturation ($0.90 \le c_t < 1.00$), high congestion ($0.75 \le c_t < 0.90$) or
moderate ($c_t < 0.75$). Figure~\ref{fig:trigger_reasons_parity} shows each trigger's
share of the attack arm's attacked steps. The random control passes through the same
gate, and its composition is nearly identical: at $b = 0.15$, for instance, its tie
share differs from the attack arm's by at most three percentage points for every
variant.

\begin{figure}[htbp]
  \centering
  \includegraphics[width=\textwidth]{trigger_reasons_independent_parity.png}
  \caption{Independent scope, FGSM-parity domain clamp: why the timing gate fired in
           the attack arm ($\varepsilon = 0.30$, $n = 30$ episodes). One panel per
           trigger that occurred in at least one run, combining the case of
           Equation~\eqref{eq:timing_rule} that admitted the step with the
           congestion band of its score. Rows are victim variants and columns the
           requested timing budget. Each cell gives the trigger's share of that run's
           attacked steps, with the count beneath it, so the three panels of one cell
           sum to 100\%. Ungated runs have no gate and are omitted.}
  \label{fig:trigger_reasons_parity}
\end{figure}

Three findings stand out. First, no attacked step in any run, in either arm, scored
$1.00$ or more, or less than $0.75$. The point mass in the score distribution that the
rationing rule of Section~\ref{subsec:timed_attack} exists to handle lies at $0.90$
rather than at full saturation. At $b = 0.10$ and $b = 0.15$ the mean threshold at
which the gate fired is $0.90$ for every variant, and the tied steps are the steps
sitting on that point mass. Its stored value falls a floating-point rounding error below
$0.90$, which is why ties are recorded in the high-congestion band although they sit at
that band's upper edge.

Second, the most congested states always fire. Steps scoring between $0.90$ and $1.00$
are admitted at every budget, in a nearly constant number per variant: $321$ for
\texttt{CC-Simple} and between $505$ and $599$ for the others. Their share of the
attacks therefore falls from $42$--$77\%$ at $b = 0.10$ to $17$--$32\%$ at $b = 0.25$.

Third, the route by which the rest of the budget is spent changes with the budget. At
$b = 0.10$ and $b = 0.15$ the threshold sits on the point mass, and ties admitted by
rationing account for $23$--$58\%$ and $45$--$66\%$ of the attacks respectively. On the
observed score sequences, admitting every tie would have raised the attack rate to
between $17\%$ and $21\%$ at both budgets, collapsing the two tightest budgets into
nearly the same experiment; rationing is what keeps them apart. From $b = 0.20$ the
threshold begins to fall below the point mass, and by $b = 0.25$ it lies between $0.82$
and $0.90$. The steps on the point mass then pass as strictly above the threshold, and
above-threshold steps in the high-congestion band, the point mass among them, make up
$57$--$77\%$ of the attacks at $b = 0.25$, while ties fall to $6$--$11\%$.

% ─────────────────────────────────────────────────────────────────────────────
\section{Results Under the Full Domain Clamp}
\label{sec:results_full_clamp}

Every result in this section, for both scopes and for both the attack and the random
control, was obtained under the \emph{full domain clamp}: after the $\ell_\infty$
projection, every feature of the perturbed observation is re-projected onto $[0,1]$.
The clamp is recorded in each evaluation record, and it is the same for all 70
configurations reported here. As set out in Section~\ref{subsec:learned_adversary},
this is stricter than the domain constraint of the FGSM study of
Section~\ref{subsec:base_fgsm}, which clamped only the normalised bandwidth features.
The admissible set used here is therefore a strict subset of the one used by the
baseline, with roughly a third of the available perturbation range removed.

Three consequences follow. First, because the attack and its random control share the
same set, the adversarial gap remains internally valid, and every comparison
\emph{within} this section is sound, including the comparison between scopes. Second,
the damage measured here is not a lower bound on the damage attained under the
baseline's admissible set, even though that set is larger: the adversaries of
Section~\ref{sec:results_parity_clamp} were trained separately, and
Section~\ref{sec:results_clamp_effect} shows them doing less damage than the adversaries
of this section on several variants. Third, a direct comparison of absolute PDR figures
against the published FGSM numbers is not licensed, since the attacker here operates
under a different constraint; that comparison belongs to the FGSM-parity clamp.

% ─────────────────────────────────────────────────────────────────────────────
\subsection{Independent scope: sensitivity to the timing budget}
\label{subsec:results_ind_timing}

\begin{figure}[htbp]
  \centering
  \includegraphics[width=\textwidth]{timing_sweep_gap_independent.png}
  \caption{Independent scope, full domain clamp. Adversarial gap for all seven victim
           variants at $\varepsilon = 0.30$, $n = 30$ paired episodes.
           \textbf{Left:} gap against the requested timing budget, with the mean
           realised attack rate printed beneath each tick and the ungated run
           included as the 100\% reference; filled markers denote a 95\% confidence
           interval excluding zero, hollow markers a null. \textbf{Right:} the same
           gaps plotted against the fraction of routing decisions the attack
           actually changed, on a logarithmic axis, with each variant drawn as one
           path walking from $b = 0.10$ through to the ungated run.}
  \label{fig:sweep_independent}
\end{figure}

The left panel of Figure~\ref{fig:sweep_independent} shows the gap rising
monotonically with the timing budget for every variant that is attackable at all.
The realised attack rates printed under the ticks confirm that the gate spends the
budget it is given: requested budgets of $0.10$, $0.15$, $0.20$ and $0.25$ realise,
averaged over the seven variants, as $10.0\%$, $15.0\%$, $19.9\%$ and $24.9\%$
respectively. This matters because the critical-state score has a point mass, which
Section~\ref{subsec:results_parity_triggers} locates at $0.90$, so a large block of
timesteps ties at exactly the trigger threshold at the tightest budgets; without the
tie rationing described in Section~\ref{subsec:timed_attack} the realised rate cannot
fall below the fraction of steps on that point mass, and the two tightest budgets
collapse into nearly the same experiment.

Table~\ref{tab:gap_sweep_independent} gives the underlying values. The architectural
pattern is sharp: the duelling value head is the vulnerable axis. All three duelling
variants without a local-critic GNN reach a significant gap at full rate, led by
\texttt{CC-Duelling-GNN} at $+2.81$~pp and \texttt{LC-Duelling} at $+2.27$~pp, while
every simple-head variant remains at or below $+0.22$~pp. Three variants,
\texttt{CC-Simple-GNN}, \texttt{LC-Simple} and \texttt{LC-Duelling-GNN}, never
separate from zero at any budget.

\begin{table}[htbp]
  \centering
  \caption{Independent scope, full domain clamp. Adversarial gap in percentage points
           by victim variant and requested timing budget, at $\varepsilon = 0.30$ and
           $n = 30$ paired episodes. Bold entries have a 95\% confidence interval
           excluding zero. Positive values indicate the attack cost the victim more
           than the random control of the same magnitude; negative values indicate
           the reverse.}
  \label{tab:gap_sweep_independent}
  \begin{tabular}{lrrrrr}
    \hline
    \textbf{Variant} & \textbf{0.10} & \textbf{0.15} & \textbf{0.20} & \textbf{0.25} & \textbf{ungated} \\
    \hline
    CC-Simple        & $-0.01$ & $+0.00$ & $+0.03$ & $\mathbf{+0.07}$ & $\mathbf{+0.22}$ \\
    CC-Duelling      & $+0.13$ & $+0.17$ & $+0.19$ & $+0.24$ & $\mathbf{+1.26}$ \\
    CC-Simple-GNN    & $-0.01$ & $+0.00$ & $+0.00$ & $-0.01$ & $-0.04$ \\
    CC-Duelling-GNN  & $\mathbf{+0.75}$ & $\mathbf{+0.77}$ & $\mathbf{+1.04}$ & $\mathbf{+1.22}$ & $\mathbf{+2.81}$ \\
    LC-Simple        & $-0.12$ & $-0.32$ & $-0.28$ & $-0.25$ & $-0.36$ \\
    LC-Duelling      & $\mathbf{+0.34}$ & $\mathbf{+0.55}$ & $\mathbf{+0.58}$ & $\mathbf{+0.71}$ & $\mathbf{+2.27}$ \\
    LC-Duelling-GNN  & $+0.00$ & $+0.00$ & $+0.01$ & $+0.01$ & $-0.05$ \\
    \hline
  \end{tabular}
\end{table}

Under this clamp, the budget sweep also supports the efficiency claim that motivated
the timing gate.
Normalising each gap by the realised attack rate gives the damage bought per
attacked timestep, and for \texttt{CC-Duelling-GNN} this falls from $7.3$~pp per
unit rate at $b = 0.10$ to $2.8$~pp when ungated, a factor of $2.6$; for
\texttt{LC-Duelling} the same ratio is $1.5$. For these adversaries, concentrating the
budget on high-utilisation states is therefore not merely cheaper but more efficient
per step spent, which is the property that motivates dividing the attack across time
in the manner of PGD. The effect is not universal: \texttt{CC-Duelling} shows
essentially no efficiency gain, its damage scaling in proportion to the number of steps
attacked, and under the FGSM-parity clamp no variant shows one
(Section~\ref{subsec:results_parity_timing}).

One caveat applies to the ungated column. Because the adversaries were trained with
the gate at $b = 0.25$, evaluating them without a gate asks them to act on the
remaining $75\%$ of states they never encountered during training. That the ungated
run is nonetheless the strongest point for every attackable variant indicates the
learned perturbation direction generalises beyond the state distribution it was
fitted on, rather than being specialised to high-utilisation moments.

% ─────────────────────────────────────────────────────────────────────────────
\subsection{Independent scope: decision flips versus delivered damage}
\label{subsec:results_ind_flips}

The right panel of Figure~\ref{fig:sweep_independent} addresses the assumption
implicit in reporting flip rates as an attack outcome. Within any single variant
the relationship holds in the expected direction: as the budget loosens, more
routing decisions change and more damage follows, so each path slopes upward.
Across variants the relationship fails completely. At the right-hand edge of the
panel four variants sit within a narrow band of $14$--$17\%$ of decisions flipped,
yet their gaps span from $-0.36$~pp to $+2.81$~pp, a spread of more than three
percentage points at essentially identical flip rates.
Table~\ref{tab:ungated_detail_independent} gives these figures directly.

\begin{table}[htbp]
  \centering
  \caption{Independent scope, full domain clamp. Ungated results at
           $\varepsilon = 0.30$, $n = 30$ paired episodes: adversarial gap with its
           95\% confidence interval, and the fraction of routing decisions changed by
           the attack and by the random control.}
  \label{tab:ungated_detail_independent}
  \begin{tabular}{lrcrr}
    \hline
    \textbf{Variant} & \textbf{Gap (pp)} & \textbf{95\% CI} & \textbf{Attack flips} & \textbf{Random flips} \\
    \hline
    CC-Simple        & $+0.22$ & $[+0.04, +0.41]$ & $9.15\%$  & $6.80\%$  \\
    CC-Duelling      & $+1.26$ & $[+0.73, +1.78]$ & $14.14\%$ & $12.02\%$ \\
    CC-Simple-GNN    & $-0.04$ & $[-0.10, +0.02]$ & $0.50\%$  & $0.34\%$  \\
    CC-Duelling-GNN  & $+2.81$ & $[+2.03, +3.59]$ & $16.69\%$ & $15.45\%$ \\
    LC-Simple        & $-0.36$ & $[-1.06, +0.34]$ & $15.94\%$ & $12.91\%$ \\
    LC-Duelling      & $+2.27$ & $[+1.60, +2.94]$ & $16.16\%$ & $11.98\%$ \\
    LC-Duelling-GNN  & $-0.05$ & $[-0.18, +0.07]$ & $3.82\%$  & $2.18\%$  \\
    \hline
  \end{tabular}
\end{table}

\texttt{LC-Simple} is the clearest counterexample. Its path in the right panel runs
downward as it moves right, so the more routing decisions the attack changes, the
\emph{better} the network performs relative to the random control. Flipping a
routing decision is evidently not the same as flipping it to a worse path, and on
this variant the learned direction is, on balance, no worse than noise. The random
control is itself instructive here: it changes $12$--$15\%$ of decisions on the
high-flip variants, which means a raw flip count says almost nothing about whether
an attacker is doing anything an untargeted disturbance would not.

The logarithmic axis separates two failure modes that a linear axis renders
identical. \texttt{CC-Simple-GNN} never moves the argmax at all, flipping between
$0.02\%$ and $0.50\%$ of decisions across the entire sweep; the perturbation is not
being absorbed downstream, it is failing to change the selected path in the first
place. This is the signature either of a policy whose decision boundaries lie beyond
the $\varepsilon$-ball, or of an adversary that has failed to learn a perturbation able
to cross them. Since the learned adversary never differentiates the victim, gradient
masking in the usual sense cannot be the cause.
\texttt{LC-Duelling-GNN} behaves quite differently: it flips $3.82\%$ of decisions
at full rate, more than seven times \texttt{CC-Simple-GNN}'s rate, and still returns
a gap of $-0.05$~pp. Here the routing genuinely changes and the packets arrive
regardless, which is absorption in the network rather than insensitivity in the
policy. The distinction matters for what would follow: the first case might yield to
a stronger optimiser or a larger budget, whereas the second reflects path redundancy
that no observation-space perturbation of this form is likely to overcome.

% ─────────────────────────────────────────────────────────────────────────────
\subsection{Coordinated scope}
\label{subsec:results_coord}

\begin{figure}[htbp]
  \centering
  \includegraphics[width=\textwidth]{timing_sweep_gap_coordinated.png}
  \caption{Coordinated scope, full domain clamp. Adversarial gap for all seven victim
           variants at $\varepsilon = 0.30$, $n = 30$ paired episodes, with the
           perturbation emitted jointly for the whole compromised group. Panels and
           markers as in Figure~\ref{fig:sweep_independent}. A thin line joins a
           label to its series where the label had to be moved clear of the data.}
  \label{fig:sweep_coordinated}
\end{figure}

Figure~\ref{fig:sweep_coordinated} repeats the sweep with the perturbation coordinated
across the compromised agents, and Table~\ref{tab:gap_sweep_coordinated} gives the
values. The timing gate behaves identically in both scopes, realising $10.0\%$,
$15.0\%$, $19.9\%$ and $24.9\%$ of steps at the four budgets, so any difference from the
independent scope is due to the coordination of the perturbation rather than to how
often it is applied.

The architectural pattern, however, changes markedly. The two local-critic variants
without GNN encoding become the most vulnerable, and both are attacked significantly
at \emph{every} budget. \texttt{LC-Simple}, which the independent attack never harmed,
rises monotonically from $+0.63$~pp at $b = 0.10$ to $+1.08$~pp at $b = 0.25$ and
$+2.94$~pp ungated. \texttt{LC-Duelling} follows the same course, from $+0.50$~pp to
$+1.17$~pp and $+2.94$~pp. The central-critic duelling variants, in contrast, are
attacked significantly only when the adversary acts on every step:
\texttt{CC-Duelling} reaches $+2.29$~pp and \texttt{CC-Duelling-GNN} $+1.46$~pp
ungated, but neither separates from zero at any gated budget. \texttt{CC-Simple},
\texttt{CC-Simple-GNN} and \texttt{LC-Duelling-GNN} remain unaffected throughout.

\begin{table}[htbp]
  \centering
  \caption{Coordinated scope, full domain clamp. Adversarial gap in percentage points
           by victim variant and requested timing budget, at $\varepsilon = 0.30$ and
           $n = 30$ paired episodes. Bold entries have a 95\% confidence interval
           excluding zero.}
  \label{tab:gap_sweep_coordinated}
  \begin{tabular}{lrrrrr}
    \hline
    \textbf{Variant} & \textbf{0.10} & \textbf{0.15} & \textbf{0.20} & \textbf{0.25} & \textbf{ungated} \\
    \hline
    CC-Simple        & $-0.01$ & $+0.01$ & $+0.04$ & $+0.06$ & $-0.10$ \\
    CC-Duelling      & $-0.11$ & $+0.05$ & $+0.02$ & $+0.26$ & $\mathbf{+2.29}$ \\
    CC-Simple-GNN    & $+0.00$ & $+0.01$ & $+0.01$ & $+0.01$ & $+0.02$ \\
    CC-Duelling-GNN  & $+0.07$ & $-0.16$ & $-0.14$ & $-0.09$ & $\mathbf{+1.46}$ \\
    LC-Simple        & $\mathbf{+0.63}$ & $\mathbf{+0.67}$ & $\mathbf{+0.80}$ & $\mathbf{+1.08}$ & $\mathbf{+2.94}$ \\
    LC-Duelling      & $\mathbf{+0.50}$ & $\mathbf{+0.80}$ & $\mathbf{+0.92}$ & $\mathbf{+1.17}$ & $\mathbf{+2.94}$ \\
    LC-Duelling-GNN  & $+0.00$ & $-0.01$ & $+0.00$ & $+0.00$ & $-0.05$ \\
    \hline
  \end{tabular}
\end{table}

Where coordination works, dividing the attack across time remains efficient. The
damage bought per attacked step for \texttt{LC-Simple} is $6.3$~pp per unit rate at
$b = 0.10$ against $2.9$~pp ungated, a factor of $2.2$, and for \texttt{LC-Duelling} it
is $5.0$~pp against $2.9$~pp, a factor of $1.7$. The right panel repeats the lesson of
Section~\ref{subsec:results_ind_flips}: the paths of \texttt{LC-Simple} and
\texttt{LC-Duelling} nearly coincide, ending at $+2.94$~pp with $14.7\%$ and $15.6\%$ of
decisions flipped, while \texttt{CC-Duelling-GNN} ends at only $+1.46$~pp despite
flipping $16.0\%$, more than either.

% ─────────────────────────────────────────────────────────────────────────────
\subsection{Effect of coordination}
\label{subsec:results_coordination_effect}

The two scopes can be compared exactly. Both were evaluated on the same traffic seed,
and because neither the clean arm nor the random control depends on the adversary,
their per-episode PDR series are bit-identical across the two scopes in all 35
configurations. The difference between the coordinated and independent adversarial
gaps is then exactly the per-episode difference between the two attack arms,
\begin{equation}
  \Delta^{\text{coord}}
  = \Delta^{\text{adv}}_{\text{coord}} - \Delta^{\text{adv}}_{\text{ind}}
  = \frac{1}{n} \sum_{e=1}^{n}
    \Big( \mathrm{PDR}^{\text{atk,ind}}_{e} - \mathrm{PDR}^{\text{atk,coord}}_{e} \Big),
  \label{eq:coordination_effect}
\end{equation}
which is paired within each episode and carries a confidence interval of the same form
as Equation~\eqref{eq:adversarial_gap}. A positive value means coordination increased
the damage done. Table~\ref{tab:coordination_effect} reports it for every
configuration.

\begin{table}[htbp]
  \centering
  \caption{Effect of coordination under the full domain clamp: coordinated minus
           independent adversarial gap, in percentage points, by victim variant and
           requested timing budget ($\varepsilon = 0.30$, $n = 30$ paired episodes).
           Positive values mean coordination increased the damage. Bold entries
           have a 95\% confidence interval on the paired difference of
           Equation~\eqref{eq:coordination_effect} excluding zero.}
  \label{tab:coordination_effect}
  \begin{tabular}{lrrrrr}
    \hline
    \textbf{Variant} & \textbf{0.10} & \textbf{0.15} & \textbf{0.20} & \textbf{0.25} & \textbf{ungated} \\
    \hline
    CC-Simple        & $+0.00$ & $+0.01$ & $+0.01$ & $-0.01$ & $\mathbf{-0.32}$ \\
    CC-Duelling      & $-0.24$ & $-0.12$ & $-0.17$ & $+0.01$ & $\mathbf{+1.03}$ \\
    CC-Simple-GNN    & $+0.01$ & $+0.01$ & $+0.01$ & $+0.01$ & $+0.06$ \\
    CC-Duelling-GNN  & $\mathbf{-0.68}$ & $\mathbf{-0.93}$ & $\mathbf{-1.18}$ & $\mathbf{-1.32}$ & $\mathbf{-1.35}$ \\
    LC-Simple        & $\mathbf{+0.75}$ & $\mathbf{+1.00}$ & $\mathbf{+1.08}$ & $\mathbf{+1.34}$ & $\mathbf{+3.29}$ \\
    LC-Duelling      & $\mathbf{+0.17}$ & $\mathbf{+0.25}$ & $\mathbf{+0.34}$ & $\mathbf{+0.46}$ & $\mathbf{+0.67}$ \\
    LC-Duelling-GNN  & $+0.00$ & $-0.01$ & $-0.01$ & $-0.01$ & $+0.00$ \\
    \hline
  \end{tabular}
\end{table}

Coordination does something, but not uniformly. On the two local-critic variants
without GNN encoding it increases damage significantly at every budget.
For \texttt{LC-Simple} the gain grows from $+0.75$~pp at $b = 0.10$ to $+3.29$~pp
ungated, turning a variant the independent attack could not harm into one of the most
damaged in the study. For \texttt{LC-Duelling}, already vulnerable, the gain grows from
$+0.17$~pp to $+0.67$~pp. On the central-critic variants the effect is mixed. It is nil
at every gated budget except for \texttt{CC-Duelling-GNN}, where coordination
significantly \emph{reduces} damage at every budget, by between $0.68$ and $1.35$~pp.
It appears only when attacking every step, where it helps \texttt{CC-Duelling}
($+1.03$~pp) and hurts \texttt{CC-Simple} ($-0.32$~pp). The two GNN variants that
resisted the independent attack are unaffected by coordination as well.

The gains do not come from changing more decisions. On \texttt{LC-Simple} at $b = 0.25$
the coordinated attack flips fewer routing decisions than the independent one, $1.91\%$
against $2.24\%$, while the gap moves from $-0.25$~pp to $+1.08$~pp; ungated, its flip
rate falls from $15.94\%$ to $14.65\%$ as the gap rises by more than three percentage
points. The converse also occurs: on \texttt{CC-Simple} coordination raises the ungated
flip rate from $9.15\%$ to $12.94\%$ while the gap falls. What coordination changes is
\emph{which} decisions change together, and only that relates to delivered damage.

The concentration of the gains on the local-critic variants is consistent with how the
two critic domains of Equation~\eqref{eq:arch_critic_domains} are trained. A local
critic conditions only on its own agent's observation and action, so the policies it
trains never receive a signal about how several agents' routing choices combine on a
shared link. A central critic conditions on the joint action, so its policies are
trained against a value model that accounts for exactly those joint consequences. A
perturbation that steers several agents onto the same link targets what the local
critic cannot see. This explanation is consistent with the pattern but is not
established by it: in particular, it does not account for why the effect disappears
for \texttt{LC-Duelling-GNN}, and it does not predict the loss on
\texttt{CC-Duelling-GNN}.

That loss calls for a specific caution. As described in
Section~\ref{subsec:learned_adversary}, the coordinated adversary acts on the joint
observation of all fourteen compromised agents, yet stores one replay transition per
step rather than one per agent, so it learns from far fewer and far higher-dimensional
samples than the independent adversary does. The large gains on the local-critic
variants, obtained under the same condition, show that the reduced sample count alone
does not explain the overall pattern. It cannot, however, be ruled out as the cause of
the deficit on \texttt{CC-Duelling-GNN} in particular, and separating an under-trained
adversary from a genuine effect of coordination on that variant would require an
analysis of the coordinated adversary's learning signal.

% ─────────────────────────────────────────────────────────────────────────────
\section{Effect of the Domain Clamp}
\label{sec:results_clamp_effect}

The two clamps can be compared on paired episodes in the independent scope. The clean
arm is bit-identical across the clamps, as noted in
Section~\ref{sec:results_parity_clamp}, but the random control is not, since each
clamp's control draws from its own admissible set. The difference between the two
adversarial gaps is nevertheless paired within each episode,
\begin{equation}
  \Delta^{\text{clamp}}
  = \Delta^{\text{adv}}_{\text{parity}} - \Delta^{\text{adv}}_{\text{full}}
  = \frac{1}{n} \sum_{e=1}^{n}
    \Big[ \big( \mathrm{PDR}^{\text{rnd,parity}}_{e} - \mathrm{PDR}^{\text{atk,parity}}_{e} \big)
        - \big( \mathrm{PDR}^{\text{rnd,full}}_{e} - \mathrm{PDR}^{\text{atk,full}}_{e} \big) \Big],
  \label{eq:clamp_effect}
\end{equation}
and carries a confidence interval of the same form as
Equation~\eqref{eq:adversarial_gap}. A positive value means the parity-clamp adversary
did more damage, relative to its own random control, than the full-clamp adversary.
Table~\ref{tab:clamp_effect} reports it for every configuration.

\begin{table}[htbp]
  \centering
  \caption{Effect of the domain clamp in the independent scope: FGSM-parity minus full
           clamp adversarial gap, in percentage points, by victim variant and requested
           timing budget ($\varepsilon = 0.30$, $n = 30$ paired episodes). Positive
           values mean the parity-clamp adversary did more damage. Bold entries have a
           95\% confidence interval on the paired difference of
           Equation~\eqref{eq:clamp_effect} excluding zero.}
  \label{tab:clamp_effect}
  \begin{tabular}{lrrrrr}
    \hline
    \textbf{Variant} & \textbf{0.10} & \textbf{0.15} & \textbf{0.20} & \textbf{0.25} & \textbf{ungated} \\
    \hline
    CC-Simple        & $+0.00$ & $+0.02$ & $-0.05$ & $-0.03$ & $\mathbf{-0.24}$ \\
    CC-Duelling      & $-0.17$ & $-0.07$ & $-0.01$ & $-0.16$ & $-0.67$ \\
    CC-Simple-GNN    & $+0.01$ & $+0.00$ & $-0.01$ & $+0.01$ & $\mathbf{+0.08}$ \\
    CC-Duelling-GNN  & $\mathbf{-0.92}$ & $\mathbf{-1.02}$ & $\mathbf{-0.93}$ & $\mathbf{-1.37}$ & $\mathbf{-2.03}$ \\
    LC-Simple        & $\mathbf{+0.52}$ & $\mathbf{+1.33}$ & $\mathbf{+1.55}$ & $\mathbf{+1.70}$ & $\mathbf{+6.26}$ \\
    LC-Duelling      & $\mathbf{-0.82}$ & $\mathbf{-1.26}$ & $\mathbf{-1.61}$ & $\mathbf{-1.53}$ & $\mathbf{-3.58}$ \\
    LC-Duelling-GNN  & $+0.00$ & $+0.01$ & $+0.00$ & $-0.01$ & $-0.13$ \\
    \hline
  \end{tabular}
\end{table}

The clamp changes the result significantly on three variants at every budget, and not
in one direction. \texttt{LC-Simple} gains between $0.52$ and $6.26$~pp under the parity
clamp, while \texttt{CC-Duelling-GNN} loses between $0.92$ and $2.03$~pp and
\texttt{LC-Duelling} between $0.82$ and $3.58$~pp. These changes come from the attack
arm rather than from the control. Ungated, the parity-clamp attack arm delivers
$6.92$~pp less than the full-clamp attack arm on \texttt{LC-Simple}, and $2.80$ and
$3.68$~pp more on \texttt{CC-Duelling-GNN} and \texttt{LC-Duelling}, while the random
control moves by less than one percentage point on all three. On \texttt{CC-Duelling}
both arms move: ungated, the parity-clamp attack arm delivers $2.49$~pp more, but its
random control also delivers $1.82$~pp more, and the net change of $-0.67$~pp is not
significant. \texttt{CC-Simple} and \texttt{CC-Simple-GNN} change significantly only
ungated, and by less than $0.25$~pp, and \texttt{LC-Duelling-GNN} does not change.

Two conclusions follow. First, each measured gap describes a trained adversary, not a
bound on what the victim can suffer. The parity set contains the full-clamp set, so an
optimal adversary could do at least as much damage under it; the trained adversaries
instead do less on two variants, and on \texttt{LC-Duelling} the sign reverses. A null
result is correspondingly weak evidence of robustness: \texttt{LC-Simple}, which the
full-clamp adversary does not harm at any budget, is the most damaged variant under the
parity clamp, and is also harmed by the coordinated full-clamp adversary of
Section~\ref{subsec:results_coord}. Second, the architectural conclusion of
Section~\ref{subsec:results_ind_timing}, that the duelling value head is the vulnerable
axis, does not survive retraining the adversary under the baseline's admissible set.
The findings that hold under both clamps are narrower. \texttt{CC-Simple-GNN} and
\texttt{LC-Duelling-GNN} are not harmed at any budget under either clamp, and
\texttt{CC-Simple} by at most $0.22$~pp; and under both clamps the fraction of decisions
an attack changes does not predict the damage it does.

The two sets differ only outside the bandwidth block, where the parity clamp allows
features to be pushed below zero or above one, outside the $[0,1]$ range that the full
clamp enforces. Whether access to these out-of-range inputs is what makes \texttt{LC-Simple}
vulnerable, and what misleads the adversary trained against \texttt{LC-Duelling}, has
not been tested.
